\documentclass[10pt,a4paper]{amsart}
\usepackage[margin=19mm]{geometry}
\usepackage{amsmath,amssymb,mathtools}
\usepackage[scr=rsfs,cal=euler]{mathalfa}
\usepackage{xfrac}
\usepackage[unicode,hidelinks,bookmarks=false]{hyperref}
\newcommand{\ie}{i.e.\ }
\newcommand{\cf}{cf.\ }
\newcommand{\resp}{respectively\ }
\newcommand{\Subsection}{\subsection}
\providecommand{\nobreakdash}{\nobreak\hbox{-}}
\setlength{\emergencystretch}{2em}
\multlinegap0pt
\usepackage{bm}
\usepackage{bbm}
\usepackage{dsfont}
\usepackage{xspace}
\usepackage{cancel}
\usepackage{mathtools}
\usepackage{amsthm}
\makeatletter
\@ifundefined{theorem}{\newtheorem{theorem}{Theorem}}{}
\@ifundefined{lemma}{\newtheorem{lemma}[theorem]{Lemma}}{}
\makeatother
\title[On Hopf hypersurfaces of the complex hyperbolic quadric with]{On Hopf hypersurfaces of the complex hyperbolic quadric with constant principal curvatures}
\author{Haizhong Li, Hiroshi Tamaru, Zeke Yao}
\date{Source: arXiv:2510.12284v1 (CC BY 4.0)}
\begin{document}
\maketitle

\noindent\emph{Source and licence.} On Hopf hypersurfaces of the complex hyperbolic quadric with constant principal curvatures. Version arXiv:2510.12284v1 (CC BY 4.0), distributed under the Creative Commons Attribution 4.0 International licence, as recorded in the arXiv licence metadata for this exact version and shown on the abstract page https://arxiv.org/abs/2510.12284v1. Full rights notice: licenses/arxiv-2510.12284-CC-BY-4.0.txt.\par\smallskip

\noindent\textit{Editorial scope.} Counted unit: the Lemma whose source label is \texttt{lemma:5.3} in the article (source0.tex lines 1382--1393, source label \texttt{lemma:5.3}), delivered together with the article's own complete original proof of it (source0.tex lines 1395--1567, one \texttt{proof} environment of 173 lines). No mathematics is written, repaired or re-derived: statement and proof are the article's own text line for line. Required context retained uncounted: eqn:2.8 (lines 486--493); eqn:2.9 (lines 495--501); lemma:5.2 (lines 1335--1345). Every definition, display and label of those lines is retained unchanged. Omitted from this package and not redistributed: 1--485, 494--494, 502--1334, 1346--1381, 1394--1394, 1568--2570. Nothing inside a retained excerpt is omitted or altered: every retained body is delivered exactly as the article's lines are written, so no omission receipt of any kind accompanies this package. Citations: the retained excerpts carry no citation command at all, so no bibliography entry of the article is transcribed and no reference is redistributed. Reference adaptation: none was needed and none was made; every reference inside the delivered unit resolves inside it, so no unresolved reference or citation is produced. Printed slips kept literal: no printed word, symbol, hypothesis or number of the counted statement or of its proof was altered. Layout and class: the article's own class is \texttt{amsart} with packages including mathrsfs, amsfonts, amsmath, amssymb, amscd, amsbsy, bm, amsthm, enumerate, babel, hyperref, geometry, xcolor; the wrapper is the lane's pinned \texttt{amsart} editorial wrapper at 10pt on A4, which restates the theorem-like environments the retained text needs and adds the article's own macro definitions that the retained text needs (none), with extra wrapper packages (bm, bbm, dsfont, xspace, cancel, mathtools). The delivered numbering is the unit's own. Assets: no retained span contains a figure environment, a graphics-inclusion command, a PDF-page inclusion, a table or a TikZ picture; no image file is copied into this package, the package declares no asset dependency and no image file is redistributed with it. Licence: version arXiv:2510.12284v1 of this article is distributed under the Creative Commons Attribution 4.0 International licence, as recorded in the arXiv licence metadata for this exact version and shown on the abstract page https://arxiv.org/abs/2510.12284v1. A full rights notice is delivered at licenses/arxiv-2510.12284-CC-BY-4.0.txt. Characters: every retained excerpt is pure ASCII, so the delivered engine, XeLaTeX with its default Latin Modern text font, reports no missing character.
\begin{equation}\label{eqn:2.8}
\begin{split}
R(X,Y)Z=&-g(Y,Z)X+g(X,Z)Y-g(\phi Y,Z)\phi X+g(\phi X,Z)\phi Y\\
 &+2g(\phi X,Y)\phi Z-g(AY,Z)(AX)^\top+g(AX,Z)(AY)^\top\\
 &-g(JAY,Z)(JAX)^\top+g(JAX,Z)(JAY)^\top\\
 &+g(SZ,Y)SX-g(SZ,X)SY,
 \end{split}
\end{equation}

\begin{equation}\label{eqn:2.9}
\begin{split}
(\nabla_X S)Y-(\nabla_Y S)X=&-\eta(X)\phi Y+\eta(Y)\phi X+2g(\phi X,Y)\xi\\
&-g(X,AN)(AY)^\top+g(Y,AN)(AX)^\top\\
 &\qquad -g(X,A\xi)(JAY)^\top +g(Y,A\xi)(JAX)^\top,
\end{split}
\end{equation}

\begin{lemma}\label{lemma:5.2}
Let $M$ be a Hopf hypersurface of $Q^{m*}$ ($m\geq3$) with constant principal curvatures and $\mathfrak{A}$-isotropic unit normal vector field $N$. For all $\lambda,\mu$ in $\sigma(\mathcal{Q})$, we have
%
\begin{enumerate}
\item[(1)]
$\nabla_X Y+\lambda g(\phi X,Y)\xi-\lambda g(AX,Y)AN+\lambda g(AX,JY)A\xi\in V_\lambda$ for all $X,Y\in V_\lambda$;

\item[(2)]
$\nabla_X Y\bot V_\lambda$ if $X\in V_\lambda$, $Y\in V_\mu$, $\lambda\neq\mu$.
\end{enumerate}
\end{lemma}

\begin{lemma}\label{lemma:5.3}
Let $M$ be a Hopf hypersurface of $Q^{m*}$ ($m\geq3$) with constant principal curvatures and $\mathfrak{A}$-isotropic unit normal vector field $N$. Let $X\in \mathcal{Q}$ be a unit principal vector at a point $p$ with
associated principal curvature $\lambda$. For any principal orthonormal basis $\{e_i\}_{i=1}^{2m-4}$ of
$\mathcal{Q}$ satisfying $Se_i=\mu_ie_i$, we have
%
\begin{equation}\label{eqn:ca2}
\begin{aligned}
\sum_{i=1, \mu_i\neq\lambda}^{2m-4}\frac{\lambda\mu_i-1}{\lambda-\mu_i}\Big\{&1+2g(\phi X,e_i)^2-2g(AX,e_i)^2-2g(AX,Je_i)^2\\
&+g(AX,X)g(Ae_i,e_i)+g(AX,JX)g(Ae_i,Je_i)\Big\}=0.
\end{aligned}
\end{equation}
\end{lemma}

\begin{proof}
Let $\mu\neq\lambda$, $\mu\in\sigma(\mathcal{Q})$ and $Y\in \mathcal{Q}$ be another unit
principal vector at $p$ with corresponding principal
curvature $\mu$. Extend $X$ and $Y$ to be principal vector fields near $p$.

(1). First note that $g(\nabla_X Y,\xi)=-g(Y,\nabla_X \xi)=-g(Y,\phi SX)=-\lambda g(\phi X,Y)$. Similarly,
$$
g(\nabla_Y X,\xi)=-\mu g(\phi Y,X)=\mu g(\phi X,Y).
$$
Thus, $g([X,Y],\xi)=-(\lambda+\mu)g(\phi X,Y)$.

By $g(\nabla_X Y,AN)=-g(Y,\nabla_X (AN))=-g(Y,q(X)A\xi-\lambda AX)=\lambda g(A X,Y)$, and
$g(\nabla_Y X,AN)=\mu g(AX,Y)$, we have $g([X,Y],AN)=(\lambda-\mu)g(AX,Y)$.

By $g(\nabla_X Y,A\xi)=-g(Y,\nabla_X (A\xi))=g(Y,q(X)AN+\lambda JAX)=-\lambda g(AX,JY)$ and
$g(\nabla_Y X,A\xi)=-\mu g(AX,JY)$, we have $g([X,Y],A\xi)=-(\lambda-\mu)g(AX,JY)$.

Next, using Codazzi equation \eqref{eqn:2.9}, we compute
\begin{equation}\label{eqn:5.9}
\begin{aligned}
g((\nabla_{[X,Y]}S)X,Y)=&g((\nabla_X S)[X,Y],Y)-g([X,Y],\xi)g(\phi X,Y)\\
&-g([X,Y],AN)g(AX,Y)-g([X,Y],A\xi)g(JAX,Y)\\
=&g((\nabla_X S)[X,Y],Y)+(\lambda+\mu)g(\phi X,Y)^2\\
&-(\lambda-\mu)g(AX,Y)^2-(\lambda-\mu)g(AX,JY)^2\\
=&g((\nabla_Y S)X,[X,Y])+2g(\phi X,Y)g([X,Y],\xi)\\
&+(\lambda+\mu)g(\phi X,Y)^2-(\lambda-\mu)\{g(AX,Y)^2-g(AX,JY)^2\}\\
=&g((\nabla_Y S)X,[X,Y])-(\lambda+\mu)g(\phi X,Y)^2\\
&-(\lambda-\mu)g(AX,Y)^2-(\lambda-\mu)g(AX,JY)^2.
\end{aligned}
\end{equation}
But now,
$$
g(\nabla_X Y,(\nabla_Y S)X)=g((\lambda I-S)\nabla_Y X,\nabla_X Y),
$$
while
$$
\begin{aligned}
g(\nabla_Y X,(\nabla_Y S)X)&=g(\nabla_Y X,(\nabla_X S)Y)+2g(\phi Y,X)g(\nabla_Y X,\xi)\\
&=g((\nabla_X S)Y,\nabla_Y X)-2g(\phi X,Y)g(\nabla_Y X,\xi)\\
&=g((\mu I-S)\nabla_X Y,\nabla_Y X)-2\mu g(\phi X,Y)^2.
\end{aligned}
$$
Thus
\begin{equation}\label{eqn:5.10}
g([X,Y],(\nabla_Y S)X)=(\lambda-\mu)g(\nabla_X Y,\nabla_Y X)+2\mu g(\phi X,Y)^2.
\end{equation}
On substituting in equation \eqref{eqn:5.9} we obtain
%
\begin{equation}\label{eqn:5.1}
g((\nabla_{[X,Y]} S)X,Y)=(\lambda-\mu)\Big\{g(\nabla_X Y,\nabla_Y X)
-g(\phi X,Y)^2-g(AX,Y)^2-g(AX,JY)^2\Big\}.
\end{equation}

(2). By using the Gauss equation \eqref{eqn:2.8}, we have
%
\begin{equation}\label{eqn:5.2}
\begin{aligned}
g(R(X,Y)Y,X)=&-1+\lambda\mu-3g(\phi X,Y)^2+g(AX,Y)^2+g(AX,JY)^2\\
&-g(AX,X)g(AY,Y)-g(AX,JX)g(AY,JY).
\end{aligned}
\end{equation}

(3). Note that $g(\nabla_Y Y,X)=0$ by Lemma \ref{lemma:5.2}, and so
$$
\begin{aligned}
g(\nabla_X {\nabla_Y Y},X)&=-g(\nabla_Y Y,\nabla_X X)=-g(\nabla_Y Y,AN)g(\nabla_X X,AN)-g(\nabla_Y Y,A\xi)g(\nabla_X X,A\xi)\\
&=-\lambda\mu [g(AX,X)g(AY,Y)+g(AX,JX)g(AY,JY)].
\end{aligned}
$$
Similarly, $g(\nabla_X Y,X)=0$, so that
$$
g(\nabla_Y {\nabla_X Y},X)=-g(\nabla_X Y,\nabla_Y X).
$$
So, $g((\nabla_{[X,Y]} S)X,Y)=(\lambda-\mu)g(\nabla_{[X,Y]} X,Y)$. We then compute
$$
\begin{aligned}
g(\nabla_X{\nabla_Y Y}-\nabla_Y {\nabla_X Y}-\nabla_{[X,Y]} Y,X)=&-\lambda\mu [g(AX,X)g(AY,Y)+g(AX,JX)g(AY,JY)]\\
&+g(\nabla_X Y,\nabla_Y X)
+\frac{1}{\lambda-\mu}g((\nabla_{[X,Y]} S)X,Y),
\end{aligned}
$$
which gives
\begin{equation}\label{eqn:5.3}
\begin{aligned}
g(R(X,Y)Y,X)=&-\lambda\mu[g(AX,X)g(AY,Y)+g(AX,JX)g(AY,JY)]\\
&+g(\nabla_X Y,\nabla_Y X)+\frac{1}{\lambda-\mu}g((\nabla_{[X,Y]} S)X,Y).
\end{aligned}
\end{equation}

(4). For a unit principal vector $Z\in \mathcal{Q}$
corresponding to a principal curvature $\nu$ not equal to $\lambda$ or $\mu$. Extend $Z$ to be principal vector fields near $p$.
By using the Codazzi equation \eqref{eqn:2.9}, we compute
$$
g((\nabla_Z S)X,Y)=g((\nabla_X S)Z,Y)=g(Z,(\nabla_X S)Y)=(\mu-\nu)g(Z,\nabla_X Y).
$$
The same calculation with $X$ and $Y$ interchanged gives
$$
g((\nabla_Z S)X,Y)=(\lambda-\nu)g(Z,\nabla_Y X).
$$
Multiplying these two equations together gives
%
\begin{equation}\label{eqn:5.4}
(\lambda-\nu)(\mu-\nu)g(\nabla_X Y,Z)g(\nabla_Y X,Z)=g((\nabla_Z S)X,Y)^2.
\end{equation}

(5). Note that
$$
g(\nabla_X Y,\xi)g(\nabla_Y X,\xi)=g(\phi SX,Y)g(\phi SY,X)=\lambda\mu g(\phi X,Y)g(\phi Y,X)=-\lambda\mu g(\phi X,Y)^2,
$$
$$
\begin{aligned}
g(\nabla_X Y,AN)&g(\nabla_Y X,AN)=g(Y,\nabla_X (AN))g(X,\nabla_Y (AN))\\
&=g(Y,q(X)A\xi-\lambda AX)g(X,q(Y)A\xi-\mu AY)
=\lambda\mu g(AX,Y)^2,
\end{aligned}
$$
$$
\begin{aligned}
g(\nabla_X Y,A\xi)&g(\nabla_Y X,A\xi)=g(Y,\nabla_X (A\xi))g(X,\nabla_Y (A\xi))\\
&=g(Y,-q(X)AN-\lambda JAX)g(X,-q(Y)AN-\mu JAY)=\lambda\mu g(AX,JY)^2.
\end{aligned}
$$
Thus, in order to express $g(\nabla_X Y,\nabla_Y X)$ in terms of the orthonormal principal basis,
we need only observe that the terms omitted from the full summation of
the $g(\nabla_X Y,e_i)g(\nabla_Y X,e_i)$ (i.e., those $i$ for which $\mu_i=\lambda$ or $\mu_i=\mu$
actually vanish and make no contribution to the sum). This is easily checked using Lemma \ref{lemma:5.2}.
Then, we can have
%
\begin{equation}\label{eqn:5.5}
g(\nabla_X Y,\nabla_Y X)=\sum_{\mu_i\neq\lambda,\mu}g(\nabla_X Y,e_i)g(\nabla_Y X,e_i)
-\lambda\mu\Big\{g(\phi X,Y)^2-g(AX,Y)^2-g(AX,JY)^2\Big\}.
\end{equation}

(6). Combining equations \eqref{eqn:5.1}, \eqref{eqn:5.2} and \eqref{eqn:5.3}, we get
$$
\begin{aligned}
2g(\nabla_X Y,\nabla_Y X)=&-1+\lambda\mu-2g(\phi X,Y)^2+2g(AX,Y)^2+2g(AX,JY)^2\\
&+(\lambda\mu-1)[g(AX,X)g(AY,Y)+g(AX,JX)g(AY,JY)].
\end{aligned}
$$
Then, we use this and the result of equation \eqref{eqn:5.4} in equation \eqref{eqn:5.5}, we get
%
\begin{equation}\label{eqn:5.6}
\begin{aligned}
2\sum_{\mu_i\neq\lambda,\mu}\frac{g((\nabla_{e_i} S)X,Y)^2}{(\lambda-\mu_i)(\mu-\mu_i)}
=&(\lambda\mu-1)\Big\{1+2g(\phi X,Y)^2-2g(AX,Y)^2-2g(AX,JY)^2\\
&   +g(AX,X)g(AY,Y)+g(AX,JX)g(AY,JY)\Big\}.
\end{aligned}
\end{equation}

Now for any $j$ with $\mu_j\neq\lambda$, we have (setting $Y=e_j$ in equation \eqref{eqn:5.6}),
%
\begin{equation}\label{eqn:5.7}
\begin{aligned}
2\sum_{\mu_i\neq\lambda,\mu_j}\frac{g((\nabla_{e_i} S)e_j,X)^2}{(\lambda-\mu_i)(\lambda-\mu_j)(\mu_j-\mu_i)}
=&\frac{\lambda\mu_j-1}{\lambda-\mu_j}\Big\{1+2g(\phi X,e_j)^2-2g(AX,e_j)^2-2g(AX,Je_j)^2\\
&   +g(AX,X)g(Ae_j,e_j)+g(AX,JX)g(Ae_j,Je_j)\Big\}.
\end{aligned}
\end{equation}

Summing this over all $j$ for which $\mu_j\neq\lambda$, we have
%
\begin{equation}\label{eqn:5.8}
\begin{aligned}
2\sum_{i,j;\mu_i\neq\mu_j;\mu_i,\mu_j\neq\lambda}&\frac{g((\nabla_{e_i} S)e_j,X)^2}{(\lambda-\mu_i)(\lambda-\mu_j)(\mu_j-\mu_i)}\\
&=\sum_{\mu_j\neq\lambda}\frac{\lambda\mu_j-1}{\lambda-\mu_j}\Big\{1+2g(\phi X,e_j)^2-2g(AX,e_j)^2-2g(AX,Je_j)^2\\
&\qquad\qquad\qquad\qquad  +g(AX,X)g(Ae_j,e_j)+g(AX,JX)g(Ae_j,Je_j)\Big\}.
\end{aligned}
\end{equation}

Since the summand on the left side of equation \eqref{eqn:5.8} is skew-symmetric in $\{i,j\}$,
the value of the sum is $0$, and so the sum on the right is $0$.
\end{proof}

\end{document}
